%% Question bank chapter file - generated by bank_tools.py build (do not edit while the bank is still being built:
%% it is regenerated from the per-test data in Question Bank\data). Every question is one bankq block:
%%   \begin{bankq}[id=..., source={...}, paper=.., qno=.., type=..., subtopic={...}, answer={...}, answer_src=key|solved,
%%                 working={...}, difficulty=1-10, diff_parts={K S I C F}, diff_note={...},
%%                 group=<passage id or empty>, also_in={duplicate copies in other tests}]
%% difficulty = min(10, K+S+I+C+F-1): K concepts 1-3, S steps 1-4, I insight 0-2, C computation 0-2, F format 0-1
%%   ...question LaTeX, options, figures (\bfig / \bfigi -> img/)...
%%   \end{bankq}
%% Passages shared by several questions are bankpassage blocks; their questions carry group=<passage id>.
\documentclass[10pt]{article}
\usepackage[a4paper,margin=1.8cm]{geometry}
\usepackage[T1]{fontenc}
\usepackage{amsmath,amssymb,mathtools}
\usepackage{graphicx}
\usepackage[export]{adjustbox}
\usepackage{xcolor,needspace,keyval}
\usepackage{enumitem}
\usepackage{booktabs,array,tabularx}
\usepackage[hidelinks]{hyperref}
\setlength{\parindent}{0pt}
\setlength{\parskip}{3pt}
\makeatletter
\define@key{bq}{id}{\def\bq@id{#1}}\define@key{bq}{source}{\def\bq@src{#1}}\define@key{bq}{paper}{\def\bq@paper{#1}}
\define@key{bq}{qno}{\def\bq@qno{#1}}\define@key{bq}{type}{\def\bq@type{#1}}\define@key{bq}{subtopic}{\def\bq@sub{#1}}
\define@key{bq}{answer}{\def\bq@ans{#1}}\define@key{bq}{answer_src}{\def\bq@asrc{#1}}\define@key{bq}{working}{\def\bq@work{#1}}
\define@key{bq}{group}{\def\bq@group{#1}}\define@key{bq}{also_in}{\def\bq@also{#1}}
\define@key{bq}{difficulty}{\def\bq@diff{#1}}\define@key{bq}{diff_parts}{\def\bq@dparts{#1}}\define@key{bq}{diff_note}{\def\bq@dnote{#1}}
\newenvironment{bankq}[1][]{\def\bq@work{}\def\bq@also{}\def\bq@group{}\def\bq@diff{?}\def\bq@dparts{}\def\bq@dnote{}\setkeys{bq}{#1}%
  \par\needspace{4\baselineskip}\medskip\noindent{\small\textbf{\bq@id}\hfill\color{gray}[\bq@src, Paper \bq@paper, Q\bq@qno]
  \bq@type\ $\cdot$ \bq@sub\ $\cdot$ \textbf{D\bq@diff}}\par\nopagebreak\relax}
  {\par{\small\color{teal}Answer: \textbf{\bq@ans}\ifx\bq@work\@empty\else\ (solved: \bq@work)\fi
  \ifx\bq@also\@empty\else\par\color{gray}Also in: \bq@also\fi}\par}
\define@key{bp}{id}{\def\bp@id{#1}}\define@key{bp}{source}{\def\bp@src{#1}}\define@key{bp}{questions}{\def\bp@qs{#1}}
\newenvironment{bankpassage}[1][]{\setkeys{bp}{#1}\par\needspace{4\baselineskip}\bigskip\noindent
  {\small\textbf{Passage \bp@id}\hfill\color{gray}[\bp@src; questions \bp@qs]}\par\nopagebreak\relax\itshape}{\par}
\makeatother
\newcommand{\bfig}[2]{\par\smallskip{\centering\includegraphics[width=#1pt,max width=0.62\linewidth,max height=0.32\textheight]{img/#2}\par}\smallskip}
\newcommand{\bfigi}[2]{\includegraphics[width=#1pt,max width=\linewidth,max height=0.32\textheight,valign=c]{img/#2}}
\newcommand{\bmissing}[1]{\textit{[figure missing in the source conversion: \detokenize{#1}]}}
\newenvironment{qoptions}{\begin{itemize}[leftmargin=1.8em,itemsep=1pt,topsep=2pt,label={}]}{\end{itemize}}
\newcommand{\opt}[2]{\item[(#1)] #2}
\providecommand{\fwstart}{}\providecommand{\fwend}{}\providecommand{\qfullsplit}{}
\newcolumntype{L}[1]{>{\raggedright\arraybackslash}p{#1}}
\newcolumntype{C}[1]{>{\centering\arraybackslash}p{#1}}
\newcolumntype{R}[1]{>{\raggedleft\arraybackslash}p{#1}}
\providecommand{\sectionhdr}[1]{}\providecommand{\subsectionhdr}[1]{}


\begin{document}
{\Large\bfseries Chemical Bonding --- Question Bank}\par\smallskip
{\small Chemistry chapter C30. Questions from all converted test papers, grouped by source series. Answers: printed key where it matches the paper, otherwise solved (answer\_src=solved, with working). Exact repeats across papers appear once, with also\_in listing the other copies.}\par\medskip\hrule

\section*{Score-I papers}
\begin{bankpassage}[id=SI12-P2-CHEM-S1-G5-6, source={Score-I - Test 12 (09 Mar 2022)}, questions={5-6}]
\begin{center}\small
\begin{tabular}{|p{0.8cm}|p{2.6cm}|p{0.8cm}|p{4.8cm}|p{0.8cm}|p{5.2cm}|}
\hline
 & \centering Column-I\newline Reaction & & \centering Column-II\newline Change associated with reaction & & \centering Column-III\tabularnewline
\hline
(I) & N$_2\to$ N$_2^+$ & (i) & Electron is removed from gerade BMO & (P) & Magnetic behaviour of reactant and product is different\\
(II) & N$_2\to$ N$_2^-$ & (ii) & Electron is removed from ungerade BMO & (Q) & Magnetic behaviour of reactant and product is same\\
(III) & C$_2\to$ C$_2^+$ & (iii) & Electron is added to gerade ABMO & (R) & Bond order of product is more as compare to reactant\\
(IV) & O$_2\to$ O$_2^-$ & (iv) & Electron is added to ungerade ABMO & (S) & Bond order of reactant is more as compare to product\\
\hline
\end{tabular}
\end{center}
\end{bankpassage}
\begin{bankq}[id=SI12-P2-CHEM-S1-Q05, source={Score-I - Test 12 (09 Mar 2022)}, paper=2, qno=5, type=single, subtopic={MOT}, answer={C}, answer_src=key, difficulty=2, diff_parts={K1 S2 I0 C0 F0}, diff_note={concepts: molecular orbital configuration and bond order and magnetism | steps: configuration of each species; bond order; magnetic nature}, group=SI12-P2-CHEM-S1-G5-6]
Select CORRECT combination.
\begin{qoptions}
\opt{A}{(I), (ii), P}
\opt{B}{(I), (i), Q}
\opt{C}{(II), (iii), P}
\opt{D}{(II), (iv), R}
\end{qoptions}
\end{bankq}
\begin{bankq}[id=SI12-P2-CHEM-S1-Q06, source={Score-I - Test 12 (09 Mar 2022)}, paper=2, qno=6, type=single, subtopic={MOT}, answer={D}, answer_src=key, difficulty=2, diff_parts={K1 S2 I0 C0 F0}, diff_note={concepts: molecular orbital configuration and bond order and magnetism | steps: configuration of each species; bond order; magnetic nature}, group=SI12-P2-CHEM-S1-G5-6]
Select INCORRECT combination.
\begin{qoptions}
\opt{A}{(III), (ii), P}
\opt{B}{(III), (ii), S}
\opt{C}{(IV), (iii), S}
\opt{D}{(IV), (iv), Q}
\end{qoptions}

\noindent
\end{bankq}

\section*{TOAS papers}
\begin{bankq}[id=TOAS01-P1-CHEM-S1-Q02, source={TOAS - Test 01 (30 Jan 2022)}, paper=1, qno=2, type=single, subtopic={VSEPR}, answer={B}, answer_src=key, difficulty=1, diff_parts={K1 S1 I0 C0 F0}, diff_note={concepts: structure of dimeric halides | steps: bridging chlorine in Al2Cl6 and I2Cl6; planarity of each}]
Which of the following structure is non-planar and have atleast one `Cl' as bridging atom?
\begin{qoptions}
\opt{A}{$\mathrm{Be_2Cl_4}$}
\opt{B}{$\mathrm{Al_2Cl_6}$}
\opt{C}{$\mathrm{Si_2Cl_6}$}
\opt{D}{$\mathrm{I_2Cl_6}$}
\end{qoptions}
\end{bankq}
\begin{bankq}[id=TOAS01-P1-CHEM-S2-Q07, source={TOAS - Test 01 (30 Jan 2022)}, paper=1, qno=7, type=multi, subtopic={MOT}, answer={B,D}, answer_src=key, difficulty=3, diff_parts={K1 S2 I0 C0 F1}, diff_note={concepts: MO configurations of B2 | steps: ground state configuration of B2; Hund's rule; s-p mixing; Pauli principle}]
Consider four configurations of $\mathrm{B_2}$ and select correct statements(s)-\\
(I) $\sigma1s^2,\sigma^*1s^2,\sigma2s^2\ ,\ \sigma^*2s^2\ ,\sigma2p_x^2$\\
(II) $\sigma1s^2,\sigma^*1s^2,\ \sigma2s^2\ ,\ \sigma^*2s^2\ ,\pi2p_y^2$\\
(III) $\sigma1s^2\ ,\sigma^*1s^2,\sigma2s^2,\sigma^*2s^2,\pi2p_x^1=\pi2p_y^1$\\
(IV) $\sigma1s^3\ ,\sigma^*1s^3,\sigma2s^3,\sigma^*2s^1$
\begin{qoptions}
\opt{A}{(I) is ground state configuration}
\opt{B}{(II) violate Hund's rule}
\opt{C}{(III) violate s-p mixing}
\opt{D}{(IV) violate Pauli exclusion principle}
\end{qoptions}
\end{bankq}
\begin{bankq}[id=TOAS01-P1-CHEM-S2-Q10, source={TOAS - Test 01 (30 Jan 2022)}, paper=1, qno=10, type=multi, subtopic={MOT}, answer={A,B,D}, answer_src=key, difficulty=2, diff_parts={K1 S1 I0 C0 F1}, diff_note={concepts: MOs from 2p-2p overlap | steps: identify the four MOs; order of energies and occupancy in F2 and O2}]
Consider following molecular orbitals (MO) which are formed by LCAO of 2p-2p orbitals and select correct option
\bfig{306.6}{TOAS01_p1_chem_s2_q10_1.png}

\begin{qoptions}
\opt{A}{MO-IV is a antibonding molecular orbital}
\opt{B}{In $\mathrm{F_2}$ MO-IV is the HOMO \& MO-II is the LUMO}
\opt{C}{In case of $\mathrm{O_2}$ one unpaired electron in MO-III and other one in MO-IV respectively}
\opt{D}{Energy of molecular orbitals in $\mathrm{O_2^{2-}}$ is I $<$ III $<$ IV $<$ II}
\end{qoptions}
\end{bankq}
\begin{bankpassage}[id=TOAS01-P1-CHEM-S3-G3-4, source={TOAS - Test 01 (30 Jan 2022)}, questions={3-4}]
In the heavier elements of p-block elements, however as we expect form the ``inert pair effect, the lower oxidation state becomes more common on descending the group.'' Covalent nature in ionic compound is explained on the basis of polarisation effect.
\end{bankpassage}
\begin{bankq}[id=TOAS01-P1-CHEM-S3-Q03, source={TOAS - Test 01 (30 Jan 2022)}, paper=1, qno=3, type=numerical, subtopic={Bond properties}, answer={30}, answer_src=key, difficulty=2, diff_parts={K1 S1 I0 C0 F1}, diff_note={concepts: Fajans' rules | steps: ZnCl2 has the polarising cation with pseudo-inert-gas core; atomic number 30}, group=TOAS01-P1-CHEM-S3-G3-4]
Atomic number of metal in the compound having most covalent character among the $\mathrm{ZnCl_2}$ , KCl, $\mathrm{CaCl_2}$ and NaCl?
\end{bankq}
\begin{bankq}[id=TOAS01-P1-CHEM-S3-Q04, source={TOAS - Test 01 (30 Jan 2022)}, paper=1, qno=4, type=numerical, subtopic={Bond properties}, answer={4}, answer_src=key, difficulty=2, diff_parts={K1 S1 I0 C0 F1}, diff_note={concepts: inert pair effect | steps: most stable oxidation states of Tl, Pb, Bi, Sn; the maximum value}, group=TOAS01-P1-CHEM-S3-G3-4]
Among the $\mathrm{Tl^{+x}}$, $\mathrm{Pb^{+x}}$ , $\mathrm{Bi^{+x}}$ , $\mathrm{Sn^{+x}}$ the maximum possible value of $x$ will be.\\
($x$ is most stable oxidation state of that element)
\end{bankq}
\begin{bankq}[id=TOAS01-P1-CHEM-S4-Q09, source={TOAS - Test 01 (30 Jan 2022)}, paper=1, qno=9, type=numerical, subtopic={Bond properties}, answer={5}, answer_src=key, difficulty=3, diff_parts={K1 S2 I0 C0 F1}, diff_note={concepts: Fajans' rules | steps: AlCl3 and MgCl2 by cationic charge; BeCl2 and LiCl by size and charge; PbCl4 and PbCl2 by oxidation state; HgS/ZnS and NiCl2/CaCl2 by pseudo-inert-gas core}]
Which of the following pair having correct order of covalent character\\
($\mathrm{AlCl_3>MgCl_2}$), ($\mathrm{BeCl_2>LiCl}$), ($\mathrm{PbCl_4>PbCl_2}$), (HgS $>$ ZnS), ($\mathrm{NiCl_2>CaCl_2}$)
\end{bankq}
\begin{bankq}[id=TOAS01-P2-CHEM-S1-Q07, source={TOAS - Test 01 (30 Jan 2022)}, paper=2, qno=7, type=single, subtopic={Bond properties}, answer={B}, answer_src=key, difficulty=1, diff_parts={K1 S1 I0 C0 F0}, diff_note={concepts: pi-d pi bonding | steps: SiF4 has back donation into Si d orbitals}]
p$\pi$ - d$\pi$ bond is present in :
\begin{qoptions}
\opt{A}{$\mathrm{BF_2Cl}$}
\opt{B}{$\mathrm{SiF_4}$}
\opt{C}{CO}
\opt{D}{$\mathrm{BeCl_2}$}
\end{qoptions}
\end{bankq}
\begin{bankq}[id=TOAS01-P2-CHEM-S1-Q08, source={TOAS - Test 01 (30 Jan 2022)}, paper=2, qno=8, type=single, subtopic={Hybridisation}, answer={B}, answer_src=key, difficulty=2, diff_parts={K1 S1 I1 C0 F0}, diff_note={concepts: odd electron species | steps: identify the orbital holding the unpaired electron in each | I: unpaired electron location}]
In which of the following species odd electron present in pure orbital?
\begin{qoptions}
\opt{A}{$\mathrm{NO_2}$}
\opt{B}{$\mathrm{ClO_2}$}
\opt{C}{$\mathrm{CF_3}$}
\opt{D}{$\mathrm{ClO_3}$}
\end{qoptions}
\end{bankq}
\begin{bankq}[id=TOAS01-P2-CHEM-S1-Q09, source={TOAS - Test 01 (30 Jan 2022)}, paper=2, qno=9, type=single, subtopic={MOT}, answer={D}, answer_src=key, difficulty=1, diff_parts={K1 S1 I0 C0 F0}, diff_note={concepts: MO magnetism | steps: unpaired electrons for N2 and its ions; for O2 and its ions}]
In which of the following processes, the value of magnetic moment does not change ?
\begin{qoptions}
\opt{A}{$\mathrm{N_2\to N_2^-}$}
\opt{B}{$\mathrm{N_2\to N_2^+}$}
\opt{C}{$\mathrm{O_2\to O_2^+}$}
\opt{D}{$\mathrm{O_2^+\to O_2^-}$}
\end{qoptions}
\end{bankq}
\begin{bankq}[id=TOAS01-P2-CHEM-S1-Q10, source={TOAS - Test 01 (30 Jan 2022)}, paper=2, qno=10, type=single, subtopic={Hybridisation}, answer={B}, answer_src=key, difficulty=1, diff_parts={K1 S1 I0 C0 F0}, diff_note={concepts: bridge bond in dimeric halide | steps: hybridisation of I and of the bridging Cl}, also_in={TOAS - Test 05 (03 Apr 2022) P2 Q17}]
The nature of overlap in the bridge bond present in $\mathrm{I_2Cl_6}$ is :
\begin{qoptions}
\opt{A}{$\mathrm{sp^3}$--p--$\mathrm{sp^3}$}
\opt{B}{$\mathrm{sp^3d^2}$--$\mathrm{sp^3}$--$\mathrm{sp^3d^2}$}
\opt{C}{$\mathrm{sp^3d}$--$\mathrm{sp^3}$--$\mathrm{sp^3d}$}
\opt{D}{$\mathrm{sp^3d}$--$\mathrm{sp^2}$--$\mathrm{sp^3d}$}
\end{qoptions}
\end{bankq}
\begin{bankq}[id=TOAS01-P2-CHEM-S1-Q12, source={TOAS - Test 01 (30 Jan 2022)}, paper=2, qno=12, type=single, subtopic={Bond properties}, answer={D}, answer_src=key, difficulty=1, diff_parts={K1 S1 I0 C0 F0}, diff_note={concepts: Fajans' rules | steps: each pair has the more polarising cation second}]
Which of the following order is correct for covalent character (if the size of cations are identical in pairs)-
\begin{qoptions}
\opt{A}{$\mathrm{NaCl<CuCl}$}
\opt{B}{$\mathrm{CaCl_2<PdCl_2}$}
\opt{C}{$\mathrm{RbCl<AuCl}$}
\opt{D}{All of these}
\end{qoptions}
\end{bankq}
\begin{bankq}[id=TOAS01-P2-CHEM-S1-Q13, source={TOAS - Test 01 (30 Jan 2022)}, paper=2, qno=13, type=single, subtopic={Bond properties}, answer={D}, answer_src=key, difficulty=1, diff_parts={K1 S1 I0 C0 F0}, diff_note={concepts: trends of inorganic compounds | steps: stability and covalent character; melting point and solubility}, also_in={TOAS - Test 09 (30 May 2022) P2 Q6}]
Select \textbf{INCORRECT} order against property indicated--
\begin{qoptions}
\opt{A}{$\mathrm{AsCl_3<SbCl_3<BiCl_3}$ (relative stability)}
\opt{B}{$\mathrm{ZnO>MgO>CaO>BaO}$ (Covalent character)}
\opt{C}{$\mathrm{NaF<MgF_2<AlF_3}$ (M.P.) Melting point}
\opt{D}{$\mathrm{BeF_2>MgF_2>CaF_2>SrF_2>BaF_2}$ (water solubility)}
\end{qoptions}
\end{bankq}
\begin{bankq}[id=TOAS03-P1-CHEM-S1-Q05, source={TOAS - Test 03 (27 Feb 2022)}, paper=1, qno=5, type=single, subtopic={MOT}, answer={C}, answer_src=key, difficulty=2, diff_parts={K1 S2 I0 C0 F0}, diff_note={concepts: bond order and magnetism from MO configurations | steps: N2 to N2+; CO to CO+; O2 to O2 2-; B2 to B2+}]
Which of the following statement is correct?
\begin{qoptions}
\opt{A}{for $\mathrm{N_2\to N_2^+}$ bond order decreases and magnetic nature does not change.}
\opt{B}{for $\mathrm{CO\to CO^+}$ bond order decreases and magnetic nature changes.}
\opt{C}{for $\mathrm{O_2\to O_2^{2-}}$ spin multiplicity decreases and magnetic nature changes.}
\opt{D}{for $\mathrm{B_2\to B_2^+}$ bond order increases and magnetic nature does not changes.}
\end{qoptions}
\end{bankq}
\begin{bankq}[id=TOAS03-P1-CHEM-S1-Q06, source={TOAS - Test 03 (27 Feb 2022)}, paper=1, qno=6, type=single, subtopic={Hybridisation}, answer={B}, answer_src=key, difficulty=2, diff_parts={K1 S2 I0 C0 F0}, diff_note={concepts: hybridisation and s-character | steps: ions in each solid; hybridisation of cation and anion; difference of s character}]
In which of the following maximum difference of \% s-character between hybrid orbital of cationic and anionic part.
\begin{qoptions}
\opt{A}{$\mathrm{Cl_2O_{6(s)}}$}
\opt{B}{$\mathrm{N_2O_{5(s)}}$}
\opt{C}{$\mathrm{I_2Cl_{6(\ell)}}$}
\opt{D}{$\mathrm{PCl_{5(s)}}$}
\end{qoptions}
\end{bankq}
\begin{bankq}[id=TOAS03-P1-CHEM-S2-Q12, source={TOAS - Test 03 (27 Feb 2022)}, paper=1, qno=12, type=multi, subtopic={Bond properties}, answer={A,B,D}, answer_src=key, difficulty=3, diff_parts={K1 S2 I0 C0 F1}, diff_note={concepts: Fajans' rules | steps: Ag2S versus Ag2O; Fe(OH)3 versus Fe(OH)2; BaCO3 versus MgCO3; BeCO3 stability}]
Which of the following can be explained on the basis of Fajan's Rules.
\begin{qoptions}
\opt{A}{$\mathrm{Ag_2S}$ is less soluble than $\mathrm{Ag_2O}$.}
\opt{B}{$\mathrm{Fe(OH)_3}$ is less soluble than $\mathrm{Fe(OH)_2}$.}
\opt{C}{$\mathrm{BaCO_3}$ is less soluble than $\mathrm{MgCO_3}$.}
\opt{D}{$\mathrm{BeCO_3}$ is thermally less stable than $\mathrm{BaCO_3}$}
\end{qoptions}
\end{bankq}
\begin{bankq}[id=TOAS03-P1-CHEM-S3-Q05, source={TOAS - Test 03 (27 Feb 2022)}, paper=1, qno=5, type=numerical, subtopic={MOT}, answer={4}, answer_src=key, difficulty=3, diff_parts={K1 S2 I0 C0 F1}, diff_note={concepts: MO bond order and paramagnetism | steps: bond orders of the ions; unpaired electrons of each; count fractional and paramagnetic ones}]
Find the no. of species having fractional bond order and paramagnetic in nature?

$\mathrm{N_2^+,\ O_2,\ O_2^+,\ F_2,\ B_2^+,\ C_2^+,\ CN^-}$
\end{bankq}
\begin{bankq}[id=TOAS03-P2-CHEM-S1-Q05, source={TOAS - Test 03 (27 Feb 2022)}, paper=2, qno=5, type=multi, subtopic={Bond properties}, answer={B,C,D}, answer_src=key, difficulty=3, diff_parts={K1 S1 I1 C0 F1}, diff_note={concepts: lone pair delocalisation with silyl groups | steps: SiH3 groups delocalise the lone pair into silicon d orbitals; effect of each replacement | I: back bonding reduces donor power}]
Lone pair donating power of $\mathrm{(H_3C)N(SiH_3)_2}$ molecule is enhances by :-
\begin{qoptions}
\opt{A}{Replacing only--$\mathrm{CH_3}$ by --$\mathrm{SiH_3}$}
\opt{B}{Replacing one--$\mathrm{SiH_3}$ by --$\mathrm{CH_3}$}
\opt{C}{Replacing both--$\mathrm{SiH_3}$ by --$\mathrm{CH_3}$}
\opt{D}{Replacing both--$\mathrm{SiH_3}$ by --Hydrogen}
\end{qoptions}
\end{bankq}
\begin{bankq}[id=TOAS03-P2-CHEM-S2-Q05, source={TOAS - Test 03 (27 Feb 2022)}, paper=2, qno=5, type=numerical, subtopic={Bond properties}, answer={3}, answer_src=solved, working={p$\pi$-d$\pi$ back bonding (donor lone pair into d orbital of Si): SiH$_3$NCO, (CH$_3$)$_3$SiOH, (SiH$_3$)$_2$O}, difficulty=3, diff_parts={K1 S2 I0 C0 F1}, diff_note={concepts: p pi - d pi back bonding | steps: identify species with donor lone pair and Si d orbitals; check each of the eight; count}]
In which of the following molecule $p\pi$ -- $d\pi$ back bonding present?

$\mathrm{CHCl_3}$, $\mathrm{B_3N_3H_6}$, $\mathrm{(CH_3)_3N}$, $\mathrm{SiH_3NCO}$, $\mathrm{:CCl_2}$, $\mathrm{(CH_3)_3SiOH}$, $\mathrm{(SiH_3)_2O}$, $\mathrm{:CF_2}$
\end{bankq}
\begin{bankq}[id=TOAS03-P2-CHEM-S3-Q05, source={TOAS - Test 03 (27 Feb 2022)}, paper=2, qno=5, type=numerical, subtopic={VSEPR}, answer={1}, answer_src=key, difficulty=3, diff_parts={K1 S2 I0 C0 F1}, diff_note={concepts: shape and polarity of molecules | steps: geometry of each of the six species; planarity; polarity}]
Among the following total number of planar and polar molecules is / are

$\mathrm{Cl_2O,\ P(CH_3)_3,\ N(CH_3)_3,\ I_2Cl_6,\ {}^{\cdot}CH_3,\ NCl_3}$
\end{bankq}
\begin{bankq}[id=TOAS03-P2-CHEM-S3-Q06, source={TOAS - Test 03 (27 Feb 2022)}, paper=2, qno=6, type=numerical, subtopic={MOT}, answer={4}, answer_src=key, difficulty=3, diff_parts={K1 S2 I0 C0 F1}, diff_note={concepts: diamagnetism | steps: electron configuration of each ion; unpaired electrons; count the diamagnetic ones}]
Find the total number of compound(s), which is/are repelled by magnetic field

$\mathrm{K_2O_2,\ KO_2,\ O_2,\ NO[BF_4],\ KCN,\ Na_2C_2}$
\end{bankq}
\begin{bankq}[id=TOAS04-P1-CHEM-S2-Q05, source={TOAS - Test 04 (15 Mar 2022)}, paper=1, qno=5, type=multi, subtopic={Bond properties}, answer={B,C,D}, answer_src=key, difficulty=3, diff_parts={K1 S2 I0 C0 F1}, diff_note={concepts: planarity and back bonding | steps: boroxine ring and borazine; silyl amines and phosphines; isothiocyanates and silyl isothiocyanates; germyl amine}]
Which of the following statement(s) is/are \textbf{CORRECT} -
\begin{qoptions}
\opt{A}{$\mathrm{[B_3O_6]^{-3}}$ ion is non-planar but $\mathrm{B_3N_3H_6}$ is planar}
\opt{B}{$\mathrm{(H_3Si)_3N}$ is planar but $\mathrm{(SiH_3)_3P}$ is pyramidal}
\opt{C}{$\mathrm{H_3C}$--NCS is bent but $\mathrm{SiH_3}$-NCS is linear}
\opt{D}{$\mathrm{(CH_3)_3N}$ is pyramidal but $\mathrm{(GeH_3)_3N}$ is planar}
\end{qoptions}
\end{bankq}
\begin{bankq}[id=TOAS04-P1-CHEM-S4-Q07, source={TOAS - Test 04 (15 Mar 2022)}, paper=1, qno=7, type=numerical, subtopic={MOT}, answer={18}, answer_src=key, difficulty=3, diff_parts={K1 S2 I0 C0 F1}, diff_note={concepts: sigma molecular orbitals | steps: electron configurations of B2 and F2; MOs symmetric about the bond axis; count the electrons}]
Consider MOT for $\mathrm{B_2}$ \& $\mathrm{F_2}$ molecules \& find total number of electrons in both which is/are present in M.O. which is/are symmetrical around bond axis :-
\end{bankq}
\begin{bankq}[id=TOAS05-P1-CHEM-S1-Q06, source={TOAS - Test 05 (03 Apr 2022)}, paper=1, qno=6, type=multi, subtopic={VSEPR}, answer={B,D}, answer_src=key, difficulty=3, diff_parts={K1 S2 I0 C0 F1}, diff_note={concepts: VSEPR bond angles | steps: Se in SeOCl2; S in SF4; Sb in SbF4-}]
Which of the following given order is \textbf{CORRECT} with respect to bond angle?
\begin{qoptions}
\opt{A}{$\mathrm{Cl{-}\widehat{Se}{-}Cl}>\mathrm{Cl{-}\widehat{Se}{-}O}$ (In $\mathrm{SeOCl_2}$)}
\opt{B}{$\mathrm{F_{eq}{-}\widehat{S}{-}F_{eq}}$(in $\mathrm{SF_4}$)$>\mathrm{F_{eq}{-}\widehat{Sb}{-}F_{eq}}$ (In $\mathrm{SbF_4^{-}}$)}
\opt{C}{$\mathrm{F_{eq}{-}\widehat{S}{-}F_{eq}}$(in $\mathrm{SF_4}$)$<\mathrm{F_{eq}{-}\widehat{Sb}{-}F_{eq}}$ (In $\mathrm{SbF_4^{-}}$)}
\opt{D}{$\mathrm{Cl{-}\widehat{Se}{-}Cl}<\mathrm{Cl{-}\widehat{Se}{-}O}$ (In $\mathrm{SeOCl_2}$)}
\end{qoptions}
\end{bankq}
\begin{bankq}[id=TOAS05-P1-CHEM-S1-Q07, source={TOAS - Test 05 (03 Apr 2022)}, paper=1, qno=7, type=multi, subtopic={VSEPR}, answer={A,B,D}, answer_src=key, difficulty=3, diff_parts={K1 S2 I0 C0 F1}, diff_note={concepts: planarity and polarity | steps: shape of NO2 and ClO2; XeF5-; ClF3}]
Which of the following species is/are planar and polar?
\begin{qoptions}
\opt{A}{$\mathrm{NO_2}$}
\opt{B}{$\mathrm{ClO_2}$}
\opt{C}{$\mathrm{XeF_5^{-}}$}
\opt{D}{$\mathrm{ClF_3}$}
\end{qoptions}
\end{bankq}
\begin{bankq}[id=TOAS05-P1-CHEM-S3-Q06, source={TOAS - Test 05 (03 Apr 2022)}, paper=1, qno=6, type=numerical, subtopic={MOT}, answer={2.00}, answer_src=key, difficulty=3, diff_parts={K1 S2 I0 C0 F1}, diff_note={concepts: MO diagrams with s-p mixing | steps: configurations; unpaired electrons; count}]
If s and p-orbital mixing (s-p mixing) is operative for all species given, then how many following species is/are paramagnetic in nature ?\\
$\mathrm{O_2^-}$, $\mathrm{N_2}$, $\mathrm{C_2}$, $\mathrm{CN^-}$, $\mathrm{B_2}$, $\mathrm{O_2^{2+}}$
\end{bankq}
\begin{bankq}[id=TOAS05-P2-CHEM-S1-Q14, source={TOAS - Test 05 (03 Apr 2022)}, paper=2, qno=14, type=single, subtopic={Bond properties}, answer={C}, answer_src=key, difficulty=2, diff_parts={K1 S2 I0 C0 F0}, diff_note={concepts: bonding statements | steps: acidity of haloforms; geometry at N in germyl isocyanate; planarity of boroxine and borazine; bond angle in disiloxane}]
Which of the following statement is \textbf{CORRECT} ?
\begin{qoptions}
\opt{A}{$\mathrm{CHF_3}$ is more acidic than $\mathrm{CHCl_3}$}
\opt{B}{Geometry around `N' is linear in $\mathrm{H_3Ge{-}NCO}$}
\opt{C}{$\mathrm{[B_3O_6]^{-3}}$ \& $\mathrm{B_3N_3H_6}$ are planar}
\opt{D}{In $\mathrm{O(SiH_3)_2}$ bond angle only affected by steric repulsion}
\end{qoptions}
\end{bankq}
\begin{bankq}[id=TOAS05-P2-CHEM-S1-Q15, source={TOAS - Test 05 (03 Apr 2022)}, paper=2, qno=15, type=single, subtopic={Bond properties}, answer={C}, answer_src=key, difficulty=2, diff_parts={K1 S2 I0 C0 F0}, diff_note={concepts: bonding statements | steps: octet in NO2; hypervalent and hypovalent species; bond order in ClO3- and ClO4-}]
Which of the following statement is correct?
\begin{qoptions}
\opt{A}{Octet rule is followed by N in $\mathrm{NO_2}$}
\opt{B}{$\mathrm{BF_3}$ is hypervalent species and $\mathrm{PF_5}$ is hypovalent species}
\opt{C}{Cl--O bond order in $\mathrm{ClO_3^-}$ is less than in $\mathrm{ClO_4^-}$}
\opt{D}{All of the above}
\end{qoptions}
\end{bankq}
\begin{bankq}[id=TOAS05-P2-CHEM-S1-Q16, source={TOAS - Test 05 (03 Apr 2022)}, paper=2, qno=16, type=single, subtopic={Bond properties}, answer={A}, answer_src=key, difficulty=2, diff_parts={K1 S2 I0 C0 F0}, diff_note={concepts: periodic trends | steps: Lewis acidity order; bond angles; carbene stability; C-O bond lengths}]
Which of the following order is/are \textbf{CORRECT}?
\begin{qoptions}
\opt{A}{$\mathrm{BF_3<BCl_3<BBr_3}$ (Lewis acidic character)}
\opt{B}{$\mathrm{NO_2^+>NO_2^->NO_3^-}$ (Bond angle)}
\opt{C}{$\mathrm{{:}CH_2<{:}CF_2<{:}CCl_2}$ (Stability of Carbene)}
\opt{D}{$\mathrm{CO_3^{2-}<CO_2<CO}$ (C-O Bond length)}
\end{qoptions}
\end{bankq}
\begin{bankq}[id=TOAS05-P2-CHEM-S1-Q18, source={TOAS - Test 05 (03 Apr 2022)}, paper=2, qno=18, type=single, subtopic={MOT}, answer={A}, answer_src=key, difficulty=1, diff_parts={K1 S1 I0 C0 F0}, diff_note={concepts: pi bonds from MO configurations | steps: configuration of C2, O2, OF and NO; count pi bonds}]
According to M.O.T., which of the following molecule has maximum number of $\pi$-bond?
\begin{qoptions}
\opt{A}{$\mathrm{C_2}$}
\opt{B}{$\mathrm{O_2}$}
\opt{C}{$\mathrm{OF}$}
\opt{D}{$\mathrm{NO}$}
\end{qoptions}
\end{bankq}
\begin{bankq}[id=TOAS05-P2-CHEM-S2-Q07, source={TOAS - Test 05 (03 Apr 2022)}, paper=2, qno=7, type=numerical, subtopic={Bond properties}, answer={8}, answer_src=key, difficulty=2, diff_parts={K1 S1 I0 C0 F1}, diff_note={concepts: percentage ionic character | steps: mu100 = e d; percent = observed/mu100}]
The dipole moment of molecule $A$--$B$ is 1.6D. If bond length of $A$--$B$ molecule is 1.625\AA, then \%ionic character is calculated to be $x\%$. Calculate $x\times0.39$ ?
\end{bankq}
\begin{bankq}[id=TOAS05-P2-CHEM-S2-Q10, source={TOAS - Test 05 (03 Apr 2022)}, paper=2, qno=10, type=numerical, subtopic={MOT}, answer={4}, answer_src=solved, working={Gerade HOMO: N$_2$ ($\sigma$2p), O$_2$, O$_2^{2-}$, N$_2^{2-}$ ($\pi^*$2p); B$_2$, C$_2$, N$_2^{2+}$ have $\pi_u$ HOMO: 4}, difficulty=4, diff_parts={K1 S2 I1 C0 F1}, diff_note={concepts: MO symmetry of the HOMO | steps: configuration for each species; gerade or ungerade HOMO; count | I: sigma2p versus pi2p ordering}, also_in={TOAS - Test 09 (30 May 2022) P1 Q8}]
Total number of species in which highest occupied molecular orbital has gerade type symmetry.\\
$\mathrm{B_2}$, $\mathrm{N_2}$, $\mathrm{C_2}$, $\mathrm{N_2^{2+}}$, $\mathrm{O_2}$, $\mathrm{O_2^{2-}}$, $\mathrm{N_2^{2-}}$
\end{bankq}
\begin{bankq}[id=TOAS08-P1-CHEM-S1-Q02, source={TOAS - Test 08 (15 May 2022)}, paper=1, qno=2, type=multi, subtopic={MOT}, answer={B,C,D}, answer_src=key, difficulty=4, diff_parts={K1 S2 I1 C0 F1}, diff_note={concepts: MOT of diatomic molecules of second period | steps: configurations of Li2 to Ne2; sigma and pi occupancy; magnetic nature | I: bond type from the MO occupancy}]
Which statement(s) can be possible for homonuclear diatomic neutral molecules formed by any $2^{\text{nd}}$ period element.
\begin{qoptions}
\opt{A}{Atleast one paramagnetic species has only $\sigma$ bonds}
\opt{B}{Atleast one diamagnetic species has only $\sigma$ bonds}
\opt{C}{Atleast one diamagnetic species has only $\pi$ bonds}
\opt{D}{Atleast one paramagnetic species has only $\pi$ bonds}
\end{qoptions}
\end{bankq}
\begin{bankq}[id=TOAS08-P2-CHEM-S3-Q03, source={TOAS - Test 08 (15 May 2022)}, paper=2, qno=3, type=numerical, subtopic={Hybridisation}, answer={3.00}, answer_src=key, difficulty=2, diff_parts={K1 S1 I0 C0 F1}, diff_note={concepts: hybridisation involving d orbitals | steps: hybridisation of each of the eight species; count those using d orbitals}]
Number of specie(s) that use atleast one double dumbell shape orbital for hybridization of central atom.\newline
(i) $[\mathrm{PtCl_4}]^{2-}$\quad (ii) $\mathrm{Ni(CO)_4}$\newline
(iii) $\mathrm{PCl_5}$\quad (iv) $[\mathrm{CoF_6}]^{3-}$\newline
(v) $\mathrm{K_4[Ni(CN)_4]}$\quad (vi) $\mathrm{XeF_6}$\newline
(vii) $\mathrm{Fe(CO)_5}$\quad (viii) $\mathrm{SO_2Cl_2}$
\end{bankq}
\begin{bankq}[id=TOAS09-P2-CHEM-S2-Q01, source={TOAS - Test 09 (30 May 2022)}, paper=2, qno=1, type=numerical, subtopic={Bond properties}, answer={6}, answer_src=solved, working={Bond orders: N$_2^+$ 2.5<3, NO$^+$ 3>2.5, CN$^-$ 3>2.5, O$_2^-$ 1.5>O$_2^{2-}$ 1, O$_2$ 2<O$_2^+$ 2.5, B$_2$ 1<B$_2^-$ 1.5: all six statements true}, difficulty=3, diff_parts={K1 S2 I0 C0 F1}, diff_note={concepts: bond order and bond length | steps: bond order of each species and its ion; compare lengths; count the true statements}]
Find out the no. of correct statements :\par
(a) Bond length N$_2^+>$ Bond length N$_2$\par
(b) Bond length NO$^+<$ Bond length of NO\par
(c) Bond length CN$^-<$ Bond length of CN\par
(d) Bond length O$_2^-<$ Bond length of O$_2^{-2}$\par
(e) Bond length O$_2>$ Bond length of O$_2^+$\par
(f) Bond length B$_2>$ Bond length of B$_2^-$
\end{bankq}
\begin{bankq}[id=TOAS10-P1-CHEM-S3-Q04, source={TOAS - Test 10 (12 Jun 2022)}, paper=1, qno=4, type=numerical, subtopic={VSEPR}, answer={8.00}, answer_src=key, difficulty=5, diff_parts={K1 S3 I1 C0 F1}, diff_note={concepts: planarity of molecules | steps: geometry at B; geometry at N; methyl carbons; hydrogens that can lie in the plane; total | I: one hydrogen of each methyl lies in the plane}]
If all the bond angles w.r.t. `N' and `B' are assumed as 120$^\circ$ in BH$_2$(NMe$_2$) molecule then maximum number of atoms in one plane will be :-
\end{bankq}
\begin{bankq}[id=TOAS11-P1-CHEM-S1-Q04, source={TOAS - Test 11 (27 Jun 2022)}, paper=1, qno=4, type=single, subtopic={Bond properties}, answer={A}, answer_src=key, difficulty=1, diff_parts={K1 S1 I0 C0 F0}, diff_note={concepts: hydrogen bond strength | steps: F-H...F is strongest followed by O-H...O and N-H...N}]
Compare strength of H-bond among the following :\par
(P) H $-$ F ......... H $-$ F\par
(Q) H $-$ O ......... H $-$ O\par
(R) H $-$ N ......... H $-$ N
\begin{qoptions}
\opt{A}{P $>$ Q $>$ R}
\opt{B}{P $<$ Q $<$ R}
\opt{C}{P $>$ R $>$ Q}
\opt{D}{R $>$ P $>$ Q}
\end{qoptions}
\end{bankq}
\begin{bankq}[id=TOAS11-P2-CHEM-S1-Q08, source={TOAS - Test 11 (27 Jun 2022)}, paper=2, qno=8, type=multi, subtopic={MOT}, answer={A}, answer_src=key, difficulty=4, diff_parts={K1 S2 I1 C0 F1}, diff_note={concepts: MO diamagnetism and gerade HOMO | steps: configuration for each species; unpaired electrons; symmetry of the HOMO | I: sigma2p versus pi2p ordering}]
Which of the following is/are diamagnetic and HOMO is gerade?
\begin{qoptions}
\opt{A}{N$_2$}
\opt{B}{O$_2$}
\opt{C}{O$_2^-$}
\opt{D}{C$_2$}
\end{qoptions}
\end{bankq}
\begin{bankq}[id=TOAS12-P1-CHEM-S1-Q04, source={TOAS - Test 12 (10 Jul 2022)}, paper=1, qno=4, type=multi, subtopic={VSEPR}, answer={C,D}, answer_src=key, difficulty=2, diff_parts={K1 S1 I0 C0 F1}, diff_note={concepts: geometry from the figure | steps: bond angles at each centre; comparisons among x, y, z}]
Which of the following relationship is(are) correct for the following figure.
\bfig{81.5}{TOAS12_p1_chem_s1_q04_1.png}
\begin{qoptions}
\opt{A}{$x=y=z$}
\opt{B}{$x=y>z$}
\opt{C}{$z>120^\circ$}
\opt{D}{$x>y$}
\end{qoptions}
\end{bankq}
\begin{bankq}[id=TOAS12-P2-CHEM-S1-Q06, source={TOAS - Test 12 (10 Jul 2022)}, paper=2, qno=6, type=single, subtopic={MOT}, answer={D}, answer_src=key, difficulty=2, diff_parts={K1 S2 I0 C0 F0}, diff_note={concepts: bond order changes and oxidation | steps: bond order increase means loss from antibonding MO; decrease means loss from bonding MO; oxidation and reduction}]
Assuming that during the formation of XY$^{\oplus}$ from XY the bond order increase where as during the formation of Z$_2^{\oplus}$ from Z$_2$ the bond order decrease.\par
Based on above information which of the following statement is incorrect?
\begin{qoptions}
\opt{A}{During the formation of XY$^+$ from XY the electron is lost from antibonding molecular orbital}
\opt{B}{During the formation of Z$_2^{\oplus}$ from Z$_2$ the electron is lost from bonding molecular orbital}
\opt{C}{In above change XY undergoes in oxidation}
\opt{D}{In above change Z$_2$ undergoes in reduction}
\end{qoptions}
\end{bankq}
\begin{bankq}[id=TOAS18-P1-CHEM-S3-Q01, source={TOAS - Test 18 (30 Oct 2022)}, paper=1, qno=1, type=numerical, subtopic={Bond properties}, answer={44}, answer_src=solved, working={$\mathrm{Cr_2O_7^{2-}}$: 8 sigma bonds, 4 pi bonds, lone pairs $2\times3+4\times2+2=16$, i.e. 32 electrons: $8+4+32=44$}, difficulty=4, diff_parts={K1 S2 I1 C0 F1}, diff_note={concepts: bonding in dichromate | steps: sigma bonds; pi bonds; lone pairs on each oxygen; add}]
For K$_2$Cr$_2$O$_7$,
\par\noindent Total number of $\sigma$-bond $=x$\par\noindent Total number of $\pi$-bond $=y$\par\noindent Total number of lone pair electron $=a$\par\noindent Find the value of $(x+y+a)$
\end{bankq}
\begin{bankq}[id=TOAS18-P2-CHEM-S1-Q01, source={TOAS - Test 18 (30 Oct 2022)}, paper=2, qno=1, type=single, subtopic={MOT}, answer={D}, answer_src=key, difficulty=1, diff_parts={K1 S1 I0 C0 F0}, diff_note={concepts: magnetism of complexes | steps: number of unpaired electrons in each species}]
Which of the following is not paramagnetic.
\begin{qoptions}
\opt{A}{K$_4\!\left[\overset{+2}{\mathrm{Fe}}(\mathrm{CN})_5\mathrm{O}_2\right]$}
\opt{B}{K$_2$MnO$_4$}
\opt{C}{$\overset{+}{\mathrm K}\ \mathrm O_2^-$}
\opt{D}{KMnO$_4$}
\end{qoptions}
\end{bankq}
\begin{bankq}[id=TOAS18-P2-CHEM-S3-Q01, source={TOAS - Test 18 (30 Oct 2022)}, paper=2, qno=1, type=single, subtopic={Hybridisation}, answer={6}, answer_src=key, difficulty=3, diff_parts={K1 S3 I0 C0 F0}, diff_note={concepts: structure of melamine | steps: sigma bonds; hybridisation; resonance; planarity; uses}]
Find the number of options which are correct for melamine.
\par\noindent(a) Planar (Almost)\par\noindent(b) 15 $\sigma$ bond\par\noindent(c) 3 p$\pi$-p$\pi$ bond are present\par\noindent(d) Resonance takes place\par\noindent(e) Pure sp$^3$ hybridise atoms are = zero\par\noindent(f) All atoms are sp$^2$ hybridised\par\noindent(g) It is used for making unbreakable Utencils
\end{bankq}
\begin{bankq}[id=TOAS19-P1-CHEM-S2-Q11, source={TOAS - Test 19 (05 Feb 2023)}, paper=1, qno=11, type=single, subtopic={Bond properties}, answer={D}, answer_src=key, difficulty=3, diff_parts={K1 S3 I0 C0 F0}, diff_note={concepts: molecular properties of simple molecules | steps: plane geometry; boiling points; dipole moments; geometry around carbon; find the wrong match}]
\begin{center}\footnotesize\setlength{\tabcolsep}{2pt}
\begin{tabular}{|L{0.8cm}L{1.3cm}|L{0.6cm}L{3.9cm}|}\hline
\multicolumn{2}{|L{2.1cm}|}{\centering List-I\\(Molecules)}&\multicolumn{2}{L{4.5cm}|}{\centering List-II\\(Various facts)}\\\hline
(I)&Neo--C$_5$H$_{12}$&(P)&Maximum number of same atoms in one plane $=3$\\\hline
(II)&C$_2$H$_2$Cl$_2$&(Q)&$n$-isomer has more surface area hence more BP then iso -- and neo --\\\hline
(III)&CH$_3$Cl&(R)&Cis-form has more dipole moment than trans-form\\\hline
(IV)&CH$_4$&(S)&has more boiling point than tetrafluoromethane due to more repulsive forces in latter\\\hline
&&(T)&has more dipole moment than dichloromethane and trichloromethane\\\hline
&&(U)&have regular geometry around carbon atom(s)\\\hline
\end{tabular}\end{center}
Select the incorrect match.
\begin{qoptions}
\opt{A}{I $\to$ P}
\opt{B}{II $\to$ R}
\opt{C}{III $\to$ T}
\opt{D}{IV $\to$ S}
\end{qoptions}
\end{bankq}
\begin{bankq}[id=TOAS19-P1-CHEM-S2-Q12, source={TOAS - Test 19 (05 Feb 2023)}, paper=1, qno=12, type=single, subtopic={Bond properties}, answer={C}, answer_src=key, difficulty=2, diff_parts={K1 S2 I0 C0 F0}, diff_note={concepts: molecular properties of simple molecules | steps: plane geometry; dipole moments; geometry around carbon; find the right match}]
\begin{center}\footnotesize\setlength{\tabcolsep}{2pt}
\begin{tabular}{|L{0.8cm}L{1.3cm}|L{0.6cm}L{3.9cm}|}\hline
\multicolumn{2}{|L{2.1cm}|}{\centering List-I\\(Molecules)}&\multicolumn{2}{L{4.5cm}|}{\centering List-II\\(Various facts)}\\\hline
(I)&Neo--C$_5$H$_{12}$&(P)&Maximum number of same atoms in one plane $=3$\\\hline
(II)&C$_2$H$_2$Cl$_2$&(Q)&$n$-isomer has more surface area hence more BP then iso -- and neo --\\\hline
(III)&CH$_3$Cl&(R)&Cis-form has more dipole moment than trans-form\\\hline
(IV)&CH$_4$&(S)&has more boiling point than tetrafluoromethane due to more repulsive forces in latter\\\hline
&&(T)&has more dipole moment than dichloromethane and trichloromethane\\\hline
&&(U)&have regular geometry around carbon atom(s)\\\hline
\end{tabular}\end{center}
Select the correct match.
\begin{qoptions}
\opt{A}{I $\to$ U}
\opt{B}{II $\to$ P}
\opt{C}{III $\to$ P}
\opt{D}{IV $\to$ T}
\end{qoptions}
\end{bankq}
\begin{bankq}[id=TOAS19-P1-CHEM-S3-Q03, source={TOAS - Test 19 (05 Feb 2023)}, paper=1, qno=3, type=numerical, subtopic={VSEPR}, answer={3.50}, answer_src=key, difficulty=4, diff_parts={K2 S2 I0 C0 F1}, diff_note={concepts: molecular geometry, polarity | steps: count molecules with equal bonds; count polar hydrogen-bonded molecules; ratio}]
Consider the following molecules

CO$_2$ , BF$_3$ , H$_2$O , NH$_3$ , CH$_4$ , PCl$_5$ , SF$_6$ , IF$_7$ , PCl$_3$F$_2$ , PCl$_2$F$_3$ , XeF$_5^\ominus$ , ClF$_3$ , SF$_4$

$\dfrac{\text{number of molecules having all equal bond length}}{\text{number of polar molecules showing H - bonding}}=x$

Find $x$?
\end{bankq}
\begin{bankq}[id=TOAS19-P1-CHEM-S3-Q04, source={TOAS - Test 19 (05 Feb 2023)}, paper=1, qno=4, type=numerical, subtopic={Bond properties}, answer={1.00}, answer_src=key, difficulty=5, diff_parts={K1 S3 I1 C0 F1}, diff_note={concepts: Lewis structure of sulfate | steps: bond order; formal charges; bond counts; lone pairs and bond pairs; ratio}]
Regarding lewis dot structure of sulphate ion, consider the following statements

(a) Resonance is possible with bond order $=1.5$

(b) Each atom contains formal charge

(c) The expanded structure and lewis structure have equal number of S--O linkage

(d) Magnitude of formal charge on sulphur is equal to $\pi$ bonds in most stable expanded structure

(e) Ratio of total number of lone pairs to bond pairs is 13 : 4

(f) Formal charge of sulphur here is different from formal charge of sulphur in lewis dot structure of SO$_3$ and SO$_2$Cl$_2$

$\dfrac{\text{number of correct statements}}{\text{number of wrong statements}}=x$

Find $x$
\end{bankq}
\end{document}
